\chapter{Experimental Study}

The growing accessibility of generative AI tools has prompted renewed interest in their potential role within education, particularly in the field of Technology-Integrated Communication in Education (TICE). For modern language teachers, lesson planning represents one of the most cognitively demanding aspects of professional practice, requiring the simultaneous consideration of curriculum objectives, student needs, and pedagogical strategy. AI-assisted planning offers a potential means of reducing this burden while broadening the quality of available teaching resources.\cite{ref:gitwiki} This experiment investigates that potential by observing language teachers engage in lesson planning both with and without AI assistance, with the aim of understanding AI not as a replacement for teacher expertise, but as a practical support tool for the new generation of classroom practice.

\section{Objective}

The objective of this experiment is to understand the advantages of AI in lesson planning for modern language courses. 

\section{Participants and Examiner}

The Experiment Conductor (EC) conducts the experiment in two phases, through instructions, record management, output collection, and providing the necessary tools to the participants.
The participants, individually noted as Test Participant (TP), are a set of eight modern language teachers from a middle school where this experiment is conducted. The selected subjects are French and English. The group of selected teachers teach all four levels at the middle school. These teachers have already participated in the Introductory Questionnaire.

\vspace{5mm}

\noindent As the participants are middle school teachers who teach either French or English, a topic common to both subjects is selected. The topic is BE + ING / Present Progressive. The participants are required to create a lesson plan for a 2-hour lecture. The topic is the same for both phases; however, it is not revealed to the participants at the beginning of the experiment. The topic is written on a piece of paper and folded, which is provided to the participants when the recording starts.

\section{List of Tools}

A set of tools are listed that are provided to the TP:
\begin{itemize}
    \item A stack of A4 plain white sheets.
    \item A pen.
\end{itemize}

\vspace{5mm}

\noindent The following set of tools are used by the EC for recording:
\begin{itemize}
    \item A camera, if consent is given,
    \item A timer, if consent for camera is denied,
    \item A pen,
\end{itemize}

\vspace{5mm}

\noindent Both Phase I and Phase II tools are common, except for a laptop with a running internet browser with several tabs open for the following generative AI (Gen-AI) tools:

\begin{itemize}
    \item ChatGPT
    \item Claude AI
    \item Duck AI
    \item Google Gemini
\end{itemize}

\vspace{5mm}


\noindent The TP may select any or all of the provided AI tools and select the best lesson plan aide as per output of the AI tool (may combine outputs as well).
And a screen recording application to record all interaction of the TP with the laptop.
The EC provides the laptop with open and running applications and browser tabs during Phase II to the TP.

\section{Data Collection Scope}

The list of data collection instruments, common for the full experiment, is as follows:

\begin{itemize}
    \item Consent Form, read and signed by TP.
    \item 2 $x$ Final Lesson Plan submitted by TP, that must be photographed.
    \item 2 $x$ Video Recording / Time Recording Table\footnote[1]{To record time steps as well as note emotional response/activity during interaction with Gen-AI.} (as per consent).
\end{itemize}


\vspace{5mm}

\noindent During Phase II of the experiment, these additional data are recorded, screen recording is additionally captured. The interaction with AI tools are saved as PDF as soon as Phase II is completed.
After the experiment, a Post Test Feedback\footnote[2]{Printed hand-out form to be filled by hand.} questionnaire is filled by the TP. This includes commentaries from the TP regarding the differences noted in the two lesson plans as well as experience of the extent of support provided by AI tools.



\subsection{Disclaimer regarding data collection}

Personal information such as name, age, gender, educational background and experience, video recordings with facial expressions capture are collected during this experiment. These data are required for two aspects:

\begin{enumerate}
    \item Demographical groupism for test group candidates and their relationship to AI tools usage and experience.
    \item Facial expression used to study interactions with AI, to note emotional responses to AI outputs.
\end{enumerate}

\noindent The identity and identification marking of the TP are strictly restricted to the EC and will not be revealed at any circumstances. The TPs will be named by random identifiers if specific commentary is made on their experiment. The video recordings will be used for inferences that is required for this report and will be deleted after the final review of this thesis. The video recordings or images will not be shared with anyone under any circumstances.

\section{Methodology}

The experiment is conducted in two phases,of approximately 20–30 minutes, known as Phase I and Phase II. A few steps are common for both Phase I and Phase II. However, this section is written in a textual flow format of the entire experiment for one TP. 

\subsection{Phase I}

The first phase of the experiment is as follows:

\begin{enumerate}
    \item The EC explains the confidentiality terms to the TP and provides the consent form Appendix -A to be read and signed by TP.
    \item The EC asks verbally for consent to video recording, and if not provided, then verbally states the use of timer.
    \item The EC explains the procedure to be followed to the TP:
    \begin{enumerate}
        \item A topic hidden in the folded piece of paper is provided.
        \item When EC states verbally “Start” and hands over the piece of paper, the TP receives the topic for the experiment.
        \item The TP needs to create a 2-hour lesson plan on the given topic in their respective subject, without any external aide.
        \item The lesson plan needs to be written on an A4 white plain paper, a stack of papers is made available for this purpose.
        \item The video recording or timer ends when the TP states verbally “Finished”.
        \item The lesson plan is to be submitted to the EC.
    \end{enumerate}
    \item Now, the EC prepares the experiment area to make sure all the tools are available at an arm’s reach.
    \item The video recording device or the table to record timer steps and comments must be set up near the EC.
    \item The folded piece of paper is with the EC.
    \item The EC asks verbally “Ready?” to which the TP is required to reply “Ready”.
    \item The EC starts the timer, states “Start” and hands over the folded piece of paper to the TP.
    \item The TP unfolds the paper and starts to make a lesson plan for the given topic as required.
    \item The EC must note all expressions of the TP and little moments from the timer steps that are relevant to understand the emotion of the TP during this task, if video recording is not permitted.
    \item The TP states “Finished” when the lesson plan is ready, and the EC immediately stops the timer and records the final time information.
    \item The EC collects the lesson plan, which is required to be photographed.
    \item The EC performs a systematic check to ensure all the data is collected without errors. This marks the End of Phase I, for one TP.
\end{enumerate}

\subsection{Phase II}
The second phase continues after Phase I.
\begin{enumerate}
    \item The EC provides the instructions for Phase II to the TP:
    \begin{enumerate}
        \item The same instructions from 3.a to 3.f are repeated with additional instructions, such as,
        \item A laptop is provided to the TP with open tabs of selected Gen-AI tools.\item The TP must use one or more of the AI tools to generate a lesson plan for the 2-hour lecture for the given topic. A lesson plan that is created without the help of AI is not acceptable.
        \item The TP must submit the final lesson plan in A4 sheets of paper.
        \item The screen is always recorded for the entirety of interaction with the AI tools.
    \end{enumerate}
    \item The EC continues the experiment as in Phase I. The TP interacts with AI tools, creates a lesson plan and submits the A4 sheet to the EC as in Phase I.
    \item The timer/video recording, as well as screen recording collection, PDF of interaction with AI tools and photographs of the lesson plan are verified by the EC, thus marking the end of Phase II.
\end{enumerate}

\noindent At the end of Phase II, a feedback form is provided to the TP. Additionally, comments on the differences between the two lesson plans are collected from the TP. These provide the scope for additional inferences on the experience of using AI for lesson planning directly from the participant.

\section{Expected Outcomes}

\begin{itemize}
    \item Based on experience the quality of lesson planning differs.
    \item Difference in expertise with AI may result in various levels of outcomes. 
    \item Ai may help in reducing the preparation time. 
    \item Ai may also provide better more detailed outcomes. 
    \item What AI produced may not correspond to the test participants’ goals. 
\end{itemize}